\documentclass[../main.tex]{subfiles}
\begin{document}
\chapter{Consistent Initial Conditions}
\label{chap:initial-conditions}

\section{State of the Art}
Constrained dynamics requires initial data compatible with its algebraic constraints. Projection corrects approximate positions and velocities toward the constraint manifold, while DAE formulations can impose additional consistency requirements \cite{shabana2013dynamics,ascherpetzold1998}. This focused review covers the geometric and numerical consistency requirements used in the example; it does not compare projection algorithms or DAE packages. Bibliographic scope checked 2026-10-07.

\section{Position and velocity consistency}

For the closed-chain model, the initial position $\vect{q}_0$ must satisfy
\begin{equation}
	\vect{\Phi}(\vect{q}_0)=\vect{0}.
	\label{eq:initial-position-consistency}
\end{equation}
The initial velocity must be tangent to the constraint manifold:
\begin{equation}
	J(\vect{q}_0)\dot{\vect{q}}_0=\vect{0}.
	\label{eq:initial-velocity-consistency}
\end{equation}
Thus, a configuration can satisfy all pin positions while a proposed velocity still separates one or more pins. Both checks are needed before integration.

Initial acceleration is not an independent initial datum in this acceleration-level formulation. Once $\vect{q}_0$ and $\dot{\vect{q}}_0$ are consistent, solve the equations of motion and acceleration constraint to obtain $\ddot{\vect{q}}_0$ and the reactions $\vect{\lambda}_0$.

\section{Projection of approximate initial data}

In practice, geometry may be specified approximately. Starting from a trial position $\vect{q}_k$, linearize the residual and choose a correction $\Delta\vect{q}_k$ satisfying
\begin{equation}
	J(\vect{q}_k)\Delta\vect{q}_k
	\approx -\vect{\Phi}(\vect{q}_k),
	\qquad
	\vect{q}_{k+1}=\vect{q}_k+\Delta\vect{q}_k.
	\label{eq:position-projection-step}
\end{equation}
Because the system uses redundant coordinates, the correction is not unique. A least-squares solve can select a minimum-norm correction; repeat the linearization until the position residual meets a stated tolerance or report failure to converge. Check that the Jacobian has the expected rank before interpreting this correction as a regular projection.

After obtaining a consistent position, correct a trial velocity $\dot{\vect{q}}_{\mathrm{trial}}$ by solving
\begin{equation}
	J(\vect{q}_0)\Delta\dot{\vect{q}}
	=-J(\vect{q}_0)\dot{\vect{q}}_{\mathrm{trial}},
	\qquad
	\dot{\vect{q}}_0=\dot{\vect{q}}_{\mathrm{trial}}+\Delta\dot{\vect{q}}.
	\label{eq:velocity-projection-step}
\end{equation}
Again, a least-squares solution selects one admissible correction, not a unique physical choice. If initial velocities have physical meaning, choose a correction criterion that preserves the intended motion as closely as possible.

\subsection{Initialization algorithm}

The notebook projects only the initial state. Its operation can be summarized as follows:
\begin{enumerate}
\item Form a geometric position guess from the specified pins, compute $\vect\Phi(\vect q_{\mathrm{guess}})$ and $J(\vect q_{\mathrm{guess}})$, and verify that $J$ has the expected row rank.
\item While $\|\vect\Phi(\vect q_k)\|_\infty$ exceeds the position tolerance, solve the minimum-norm least-squares correction $J(\vect q_k)\Delta\vect q_k\approx-\vect\Phi(\vect q_k)$ and update $\vect q_{k+1}=\vect q_k+\Delta\vect q_k$.
\item Stop with a reported failure if the residual does not reach tolerance within the iteration limit; do not silently pass an inconsistent pose to the integrator.
\item At the corrected position, evaluate $J$ and solve the minimum-norm velocity correction $J\Delta\dot{\vect q}=-J\dot{\vect q}_{\mathrm{trial}}$. Set $\dot{\vect q}_0=\dot{\vect q}_{\mathrm{trial}}+\Delta\dot{\vect q}$.
\item Check both $\|\vect\Phi(\vect q_0)\|_\infty$ and $\|J(\vect q_0)\dot{\vect q}_0\|_\infty$. Then solve the constrained acceleration equations once and check force-balance and acceleration-constraint residuals before starting the trajectory integration.
\end{enumerate}
For the notebook's projected initial state, the raw pin position residual is $4.441\times10^{-16}\,\mathrm m$, grounded translation error is zero, and grounded angle error is zero. The pin velocity residual is $4.718\times10^{-16}\,\mathrm{m\,s^{-1}}$; the grounded translation-rate and angular-rate residuals are $1.561\times10^{-17}\,\mathrm{m\,s^{-1}}$ and $1.648\times10^{-17}\,\mathrm{s^{-1}}$, respectively. These values verify that particular initialization; tolerances should be selected relative to units and model scales in other applications.

\section{Checks before integration}

Before starting the solver, evaluate $\|\vect{\Phi}(\vect{q}_0)\|_\infty$ and $\|J(\vect{q}_0)\dot{\vect{q}}_0\|_\infty$. Confirm the Jacobian rank expected for the selected configuration, then solve the constrained acceleration equations and check both force-balance and acceleration-constraint residuals. Tolerances should be reported with units and interpreted relative to the model scales; there is no universal projection tolerance for every mechanism.

In the five-link notebook, position and velocity projection are performed on the initial guess before integration. Projection is not applied after every time step in that example; subsequent drift is controlled by the selected Baumgarte term and monitored separately.

\section{Summary and transition}

Consistent initialization requires position closure and tangent velocity. Acceleration follows from the constrained equations. Projection offers a practical correction for approximate initial data, but its nonuniqueness and rank assumptions should be stated. With an admissible initial state prepared, the next chapter considers time integration and constraint drift.

\ifSubfilesClassLoaded{\backmatter\bookbibliography}{}
\end{document}